% Part: normal-modal-logic
% Chapter: axioms-systems
% Section: normal-logics

\documentclass[../../../include/open-logic-section]{subfiles}

\begin{document}

\olfileid[tr]{nml}{prf}{nor}

\olsection{Normal Modal Mantıklar}

Her modal !!{formula}s kümesi, bir aksiyomlar kümesinden türetilebilen
!!{formula}s olarak kolayca nitelenemez. Modal mantıkların düzgün davranmasını
isteriz. Her şeyden önce klasik önerme mantığında !!{derive} mümkün olan her
şey yine !!{derivable} olmalıdır; elbette artık !!{formula}s
\iftag{prvBox}{$\Box$\iftag{prvDiamond}{ ve~}{}}{}
\iftag{prvDiamond}{$\Diamond$}{} da içerebilir. Bunun için bir modal mantığın bütün totoloji
örneklerini içermesini ve modus ponens altında kapalı olmasını isteriz.

\begin{defn}
  Bir \emph{modal mantık}, aşağıdaki koşulları sağlayan bir modal
  !!{formula}s kümesi $\Sigma$'dır:
  \begin{enumerate}
  \item bütün totolojileri içerir;
  \item yerine koyma altında kapalıdır; yani $!A\in\Sigma$ ve
    $!D_1,\dots,!D_n$ !!{formula}s ise
    \[
    \SSubst{!A}{\subst{!D_1}{p_1}, \dots, \subst{!D_n}{p_n}} \in \Sigma;
    \]
  \item \emph{modus ponens} altında kapalıdır; yani $!A$ ile
    $!A\lif !B\in\Sigma$ ise $!B\in\Sigma$'dır.
  \end{enumerate}
\end{defn}

Modal mantıklar için bağıntısal semantiği kullanabilmek üzere, bütün modal
modellerde geçerli !!{formula}s da mantığa katılmalıdır.
$\Ax{K}$'nin\iftag{prvDiamond}{ ve~$\Dual{}$'ın}{} bütün örnekleri
!!{derivable} olduğunda ve !!a{formula}~$!A$ !!{derivable} iken $\Box !A$
için de aynısı geçerli olduğunda bu koşulun sağlandığı ortaya çıkar. Bu koşulları
sağlayan bir modal mantığa \emph{normal} denir. (Elbette normal olmayan
modal mantıklar da vardır; ancak alışılmış bağıntısal modeller onlar için
yeterli değildir.)

\begin{defn}
  Bir modal mantık $\Sigma$, şu formülleri içeriyorsa
  \begin{align*}
    \tag{\Ax{K}} & \Box(p \lif q) \lif (\Box p \lif \Box q),
    \iftag{prvDiamond}{\\
      \tag{\Dual} & \Diamond p \liff \lnot\Box\lnot p}{}
  \end{align*}
  ve \emph{zorunlulaştırma} altında kapalıysa, yani $!A\in\Sigma$ iken
  $\Box !A\in\Sigma$ ise \emph{normaldir}.
\end{defn}

Totolojik gerektirmenin \emph{bir dünyadaki doğruluğu} koruyacak kadar
``ince taneli'' olmasına karşılık, \Nec{} kuralının yalnızca \emph{bir
modeldeki doğruluğu} (dolayısıyla bir çerçevede ya da çerçeveler sınıfındaki
geçerliliği de) koruduğuna dikkat edin.

\begin{prop}\ollabel{prop:rk}
  Her normal modal mantık şu \RK{} kuralı altında kapalıdır:
  \begin{prooftree}
    \AxiomC{$!A_1 \lif (!A_2 \lif \cdots (!A_{n-1} \lif !A_n)\cdots)$}
    \RightLabel{\RK}
    \UnaryInfC{$\Box!A_1 \lif (\Box!A_2 \lif \cdots (\Box!A_{n-1}
      \lif \Box!A_n)\cdots).$}
  \end{prooftree}
\end{prop}

\begin{proof}
  $n$ üzerinde tümevarımla: $n=1$ ise kural yalnızca \Nec{} kuralıdır ve
  her normal modal mantık \Nec{} altında kapalıdır.

  Sonucun $n-1$ için geçerli olduğunu varsayıp $n$ için gösterelim.

  Varsayalım ki
  \begin{align*}
  & !A_1 \lif (!A_2 \lif \cdots (!A_{n-1} \lif !A_n)\cdots) \in \Sigma.
  \intertext{Tümevarım varsayımından}
  & \Box!A_1 \lif (\Box!A_2 \lif \cdots \Box(!A_{n-1} \lif !A_n)\cdots)
  \in \Sigma.
  \intertext{$\Sigma$ normal bir modal mantık olduğundan $\Ax{K}$'nin
    bütün örneklerini, özellikle şunu içerir:}
  & \Box(!A_{n-1} \lif !A_n) \lif (\Box!A_{n-1} \lif \Box!A_n) \in \Sigma.
  \intertext{Modus ponens ve uygun totoloji örneklerini kullanarak şunu elde
    ederiz:}
  & \Box!A_1 \lif (\Box!A_2 \lif \cdots (\Box!A_{n-1}
  \lif \Box!A_n)\cdots) \in \Sigma. 
  \end{align*}
\end{proof}

\begin{prop}\ollabel{prop:notDiamondBot}
  Her normal modal mantık $\Sigma$, $\lnot\Diamond\lfalse$ formülünü içerir.
\end{prop}

\begin{prob}
  \olref[nml][prf][nor]{prop:notDiamondBot} önermesini ispatlayın.
\end{prob}

\begin{prop}
  $!A_1,\dots,!A_n$ !!{formula}s olsun. $!A_1,\dots,!A_n$'nin bütün örneklerini içeren
  en küçük bir modal mantık $\Sigma$ vardır.
\end{prop}

\begin{proof}
  $!A_1,\dots,!A_n$ verildiğinde $\Sigma$'yı, $!A_1,\dots,!A_n$'nin bütün örneklerini
  içeren normal modal mantıkların kesişimi olarak tanımlayın.
  $\Frm[L]$, yani bütün !!{formula}s kümesi, böyle bir modal mantık olduğundan kesişim
  boş değildir.
\end{proof}

\begin{defn}
$!A_1,\dots,!A_n$'yi içeren en küçük normal modal mantığa bir
\emph{modal sistem} denir ve $\Log{K}!A_1\dots !A_n$ ile gösterilir. En
küçük normal modal mantık $\Log{K}$ ile gösterilir.
\end{defn}

\end{document}
